\chapter*{Plan na 52 tygodnie}
\addcontentsline{toc}{chapter}{Plan na 52 tygodnie}
\markboth{Plan}{Plan}

Tabela poniżej jest \textbf{generowana} z \code{data/units.tsv}. Tydzień kalendarzowy różni się od
numeru jednostki, bo po każdym poziomie jest tydzień buforowy.

% GENERATED by tools/gen.py -- do not edit.
\begin{plantable}
1 & A1 & 1 & Dźwięki i powitania & alfabet, wymowa, Sie/du\\
2 & A1 & 2 & Przedstawianie się & sein, zaimki osobowe, pytania W-\\
3 & A1 & 3 & Liczby i dane kontaktowe & czasowniki regularne, liczebniki 0-100\\
4 & A1 & 4 & Rzeczowniki i rodzajniki & rodzaj, liczba mnoga, ein/eine\\
5 & A1 & 5 & Haben, biernik i przeczenie & haben, Akkusativ, kein/nicht\\
6 & A1 & 6 & Rodzina i zawody & przymiotniki dzierżawcze\\
7 & A1 & 7 & Czas i terminy & godziny, dni, miesiące, daty\\
8 & A1 & 8 & Jedzenie i picie & möchte, restauracja, ceny\\
9 & A1 & 9 & Zakupy i ilości & czasowniki nieregularne (a→ä, e→i)\\
10 & A1 & 10 & Mieszkanie i miejsca & przyimki in/auf/an, Wo? Wohin?\\
11 & A1 & 11 & Czasowniki modalne & können, müssen, wollen, dürfen, Satzklammer\\
12 & A1 & 12 & Plan dnia i tryb rozkazujący & czasowniki rozdzielnie złożone, Imperativ\\
\rowcolor{lightgray}13 & A1 & --- & \textit{Tydzień buforowy i test poziomu A1} & powtórka, \textit{Anki}, mock-test\\
14 & A2 & 13 & Miasto i komunikacja & Dativ, przyimki z celownikiem, droga\\
15 & A2 & 14 & Celownik i zaimki osobowe & mir/dir/ihm, czasowniki z Dativ\\
16 & A2 & 15 & Perfekt z haben & Partizip II regularny\\
17 & A2 & 16 & Perfekt z sein, czasowniki nieregularne & Partizip II nieregularny\\
18 & A2 & 17 & Präteritum: sein, haben, modalne & Präteritum, zdrowie\\
19 & A2 & 18 & Zdania podrzędne: weil, dass, wenn, ob & szyk w zdaniu podrzędnym\\
20 & A2 & 19 & Porównania i pogoda & Komparativ, Superlativ\\
21 & A2 & 20 & Czasowniki zwrotne i czas wolny & sich + Akkusativ/Dativ\\
22 & A2 & 21 & Opisujemy przymiotnikami & deklinacja przymiotników I\\
23 & A2 & 22 & Biuro, e-mail i telefon & korespondencja, grzeczność\\
24 & A2 & 23 & Plany i grzeczne prośby & Futur I, würde/könnte/hätte gern\\
25 & A2 & 24 & Podsumowanie A2 i CV & powtórka, życiorys\\
\rowcolor{lightgray}26 & A2 & --- & \textit{Tydzień buforowy i test poziomu A2} & powtórka, \textit{Anki}, mock-test\\
\end{plantable}


\begin{uwaga}[Kamienie milowe]
Po \textbf{A1}: przedstawisz się, zamówisz jedzenie, zapytasz o cenę i godzinę. Po \textbf{A2}: opowiesz
o weekendzie w czasie przeszłym, napiszesz krótki e-mail do kolegi z pracy. Po \textbf{B1}: poprowadzisz
rozmowę o swoim doświadczeniu zawodowym i zrozumiesz sedno wiadomości. Po \textbf{B2}: weźmiesz udział
w spotkaniu technicznym, zaprezentujesz architekturę i napiszesz formalną dokumentację.
To cele pracy, nie gwarancja: sprawdza je dopiero egzamin (Dodatek~E).
\end{uwaga}
